% Generated draft: requires mathematical and source review.
\documentclass[11pt]{article}
\usepackage[margin=25mm]{geometry}
\usepackage{fontspec}
\setmainfont{Latin Modern Roman}
\usepackage{amsmath,amssymb,mathtools,unicode-math}
\setmathfont{Latin Modern Math}
\usepackage{enumitem,needspace,xurl,hyperref}
\setlength{\parindent}{0pt}
\setlength{\parskip}{5pt}
\setlength{\emergencystretch}{3em}
\newcommand{\N}{\mathbb{N}}
\newcommand{\Z}{\mathbb{Z}}
\newcommand{\Q}{\mathbb{Q}}
\newcommand{\R}{\mathbb{R}}
\newcommand{\C}{\mathbb{C}}
\newcommand{\F}{\mathbb{F}}
\newcommand{\A}{\mathbb{A}}
\newcommand{\PP}{\mathbb{P}}
\newcommand{\E}{\mathbb{E}}
\newcommand{\eps}{\varepsilon}
\newcommand{\dd}{\,\mathrm{d}}
\begin{document}
\section*{UnsolvedMath ID: 30005192}
\textbf{Source problem code:} catalogue entry
\subsection{Definitions and mathematical statement}
A conventional framework for interpreting the record is the following. Regard a monoid \(M\), with identity \(e\), as a category with one object. Let \(\mathsf{Corr}\) denote the bicategory whose objects are \(C^*\)-algebras, whose arrows \(A\to B\) are right Hilbert \(B\)-modules \(E\) equipped with a nondegenerate \(*\)-homomorphism \(\varphi_E:A\to\mathcal L_B(E)\), and whose \(2\)-arrows are unitary bimodule isomorphisms. Here \(\mathcal L_B(E)\) is the algebra of adjointable operators. Composition is the interior tensor product. The compact operators \(\mathcal K_B(E)\) are the closed linear span of the operators \(\theta_{\xi,\eta}(\zeta)=\xi\langle\eta,\zeta\rangle_B\). A correspondence is called proper when \(\varphi_E(A)\subseteq\mathcal K_B(E)\).

Under this convention, a homomorphism \(F:M\to\mathsf{Corr}\), interpreted as a pseudofunctor, is a product system consisting of a \(C^*\)-algebra \(A\), correspondences \(E_m:A\to A\) for every \(m\in M\), and unitary multiplication maps
\[
\mu_{m,n}:E_m\otimes_A E_n\xrightarrow{\cong}E_{mn},\qquad E_e\cong A,
\]
with the unit identities and the associativity identities
\[
\mu_{mn,p}\circ(\mu_{m,n}\otimes 1)
=\mu_{m,np}\circ(1\otimes\mu_{n,p})
\qquad(m,n,p\in M),
\]
after the canonical tensor-product associativity identifications. This fixes a multiplication-order convention for the monoid category. In this framework, a nonproper product system means that \(\varphi_{E_m}(A)\nsubseteq\mathcal K_A(E_m)\) for at least one \(m\in M\). The supplied record does not specify these conventions or explicitly define nonproperness.

\textbf{Recorded question, conditional on these conventions.} For every monoid \(M\) and every nonproper product system \(F:M\to\mathsf{Corr}\), does a bicategorical limit exist, and, whenever it exists, how can it be described? No cancellativity, Ore condition, countability, separability, fullness, or faithfulness hypothesis is supplied.

More precisely, if ``limit'' means a conical pseudolimit, seek an object \(L\) and a pseudocone \(\lambda:\Delta L\Rightarrow F\) such that, for every \(C^*\)-algebra \(D\), composition with \(\lambda\) induces an equivalence
\[
\mathsf{Corr}(D,L)\simeq\operatorname{PsCone}(\Delta D,F),
\]
pseudonatural in \(D\). Here \(\Delta D\) is the constant diagram, and the right-hand side is the category of pseudocones and modifications. In the convention above, a pseudocone has a correspondence \(H:D\to A\) and coherent unitary identifications \(H\otimes_A E_m\cong H\) for all \(m\in M\).

\textbf{Unresolved target.} The original wording asks what a homomorphism into the correspondence bicategory is and what its ``limit'' would look like. The supplied literature assessment instead discusses colimits. A conical pseudocolimit would have the different universal property
\[
\mathsf{Corr}(L,D)\simeq\operatorname{PsCone}(F,\Delta D),
\]
with cocones given by correspondences \(H:A\to D\) and coherent unitary identifications \(E_m\otimes_A H\cong H\). The supplied evidence does not resolve whether a limit, a colimit, or a lax variant is intended. Consequently, the displayed limit question is a conditional formalization of the clean statement, not a verified reconstruction of the report's precise target.
\subsection{Short English statement}
The record asks whether an arbitrary nonproper product system over a monoid has a universal bicategorical limit and how to describe it. A conventional correspondence framework is given above, but the intended universal construction remains uncertain because the supplied literature discussion concerns colimits. The supplied status is \emph{partially solved}; this draft does not independently verify that assessment.
\subsection{Sources}
\begin{itemize}
\item Local UnsolvedMath registry, supplied complete record for permanent ID \(30005192\), titled \emph{Bicategorical Limits of Nonproper Product Systems}; exact source problem code \texttt{catalogue entry}. The recorded proposed year is \(2022\); no individual proposer is supplied.
\item Supplied report citation: \emph{Non-Commutative Geometry and Cyclic Homology (2023)}, Oberwolfach Reports, DOI \texttt{10.4171/owr/2022/35}, pp.~1987--2058. Supplied source URL and literature-source label ``original Oberwolfach report'': \url{https://doi.org/10.4171/owr/2022/35}.
\item Supplied literature source: \emph{Colimits in the correspondence bicategory}, \url{https://arxiv.org/abs/1502.07771}.
\item Supplied literature source: \emph{Groupoid models for diagrams of groupoid correspondences}, \url{https://arxiv.org/abs/2203.12068}.
\item Literature assessment dated exactly \texttt{2026-08-22}: the supplied curation labels the problem partially solved and states that homomorphisms into the correspondence bicategory and their colimits are developed for groups, groupoids, product systems, and important classes of monoids such as Ore monoids, while the unrestricted monoid formulation is broader than established results. This is supplied dated curation, not independently verified literature assessment.
\item The registry labels the statement \texttt{corrected\_verified} and reports that it was checked against the official Oberwolfach report and preserves its hypotheses, notation, and target. That verification claim is supplied metadata and has not been independently checked here. The background explicitly requests independent verification of time-sensitive status and cited claims.
\end{itemize}
\end{document}
